\newpage
\phantomsection
\label{prf:preimage-upper-lower-sets}
\LRAProofFor{thm:preimage-upper-lower-sets}

\begin{remark*}[Return]
\hyperref[thm:preimage-upper-lower-sets]{Return to Theorem}
\end{remark*}

\begin{theorem*}[Preimages of Upper and Lower Sets]
Order-preserving maps pull upper sets back to upper sets and lower sets back
to lower sets; order-reversing maps swap the two.
\end{theorem*}

\begin{proof}[Professional Standard Proof]
\LRAProofBodyStart
Let $f$ be order-preserving and let $U$ be an upper set in $Q$. If
$x\in f^{-1}(U)$ and $x\leq y$, then $f(x)\in U$ and
$f(x)\leq' f(y)$. Since $U$ is upper, $f(y)\in U$, so $y\in f^{-1}(U)$.
The lower-set case is dual. If $f$ is order-reversing, $x\leq y$ gives
$f(y)\leq' f(x)$, which swaps the upper and lower arguments.
\end{proof}

\begin{proof}[Detailed Learning Proof]
\LRAProofBodyStart
Membership in a preimage means the output lies in the target set. Monotonicity
transports the domain comparison to a target comparison. The upper- or
lower-set property then keeps the new output inside the target set.
\end{proof}
\begin{remark*}[Proof structure]
\end{remark*}



\begin{dependencies}
\begin{itemize}
  \item \hyperref[thm:preimage-upper-lower-sets]{Proof target}
\end{itemize}
\end{dependencies}

\clearpage

